%
% menu_lifecycle.tex
%
% A Z specification of the Hub's agent-menu ownership, inbox delivery, and
% session-departure lifecycle -- the state machine behind bead lux-m3xr
% (PR #495, design docs/architecture/agent-menu-dispatch.md, DES-098).
%
% WHY THIS MODEL EXISTS. The design's own Section 6 named a tripwire: the
% menu-event enqueue is safe only if it can be ordered by reusing the
% client-registry read the callback leg already proves, introducing no new
% lock and no new acquisition order. Review of the implementation found the
% tripwire tripped -- three concurrency/lifecycle defects, each of the
% class the repo's z-spec policy makes model-checking mandatory for:
%
%   D1 (inbox.py offer) -- the enqueue is NOT atomic with departure. offer
%     reads the inbox queue under _inboxes_lock, RELEASES the lock, then
%     puts. If drop_session pops the queue in that window, the put lands in
%     an orphaned queue and offer returns True: a false "delivered" signal.
%     MenuEventRouter.deliver then reports "delivered" for a click no recv()
%     will ever drain -- the click is silently lost.
%
%   D2 (departure cascade wiring) -- the inbox sink is bound into the shared
%     DepartureCascade (via inbox.ensure_writer -> bind_departure_sink) and
%     so fires on every departure trigger; the HubMenuRegistry sink is wired
%     ONLY to session_cleanup.py's graceful/SDK-idle teardown leg. On the
%     lease-lapse path (reap_lapsed_leases -> _depart_lapsed -> cascade)
%     nothing prunes the menu registry, so a timer-reaped owner's bar entry
%     never leaves _by_owner. This is the connection_lease_reaping I5/I7
%     "leaving the registry is one event, however triggered" property, now
%     applied to the menu registry rather than the inbox.
%
%   D3 (operations/menus.py set_menu) -- admit/store is a TOCTOU. admit()
%     checks the caller live-and-identified and arms its inbox under
%     StoreLock, but RELEASES it before set_menus() records the bar under
%     the menu registry's own separate lock. A departure interleaving in
%     that window stores a bar for an already-departed owner, which the
%     already-run departure sweep will never revisit -- an orphaned-owner
%     menu that outlives its session.
%
% This model refines connection_lease_reaping's departure cascade: its
% Depart / _depart_lapsed operations here sweep the MENU registry and the
% INBOX as cascade legs, and this spec proves the three delivery/ownership
% invariants over that cascade. The three defects are reproduced as
% fidelity controls (menu_lifecycle_d1_buggy.tex, _d2_buggy.tex,
% _d3_buggy.tex), each of which makes the corresponding invariant's
% negation FOUND when its guard/atomicity is removed.
%
% Read first: docs/architecture/agent-menu-dispatch.md (Sections 4-7),
% docs/connection_lease_reaping.tex (the departure cascade this refines;
% I1/I5/I7), and the code: src/punt_lux/domain/hub/inbox.py (offer,
% drop_session, ensure_writer), menu_registry.py (set_menus, drop_session),
% menu_event.py (MenuEventRouter.deliver), operations/menus.py
% (MenuOperations.set_menu), operations/menu_arming.py (MenuArming.admit),
% domain/hub/hub_display.py (_depart, _depart_lapsed, drop_connection),
% domain/hub/departure_cascade.py, session_cleanup.py.
%
% Type-check:  fuzz docs/menu_lifecycle.tex
% Model-check (deadlock-freedom + structural invariant over the state space):
%   probcli docs/menu_lifecycle.tex -model_check -p DEFAULT_SETSIZE 3
% The safety goals (M1, MO) are reachability checks; see the Verification
% section for the exact -goal predicates, following the same convention as
% connection_lease_reaping.tex.
%

\documentclass[a4paper,10pt,fleqn]{article}
\usepackage[margin=1in]{geometry}
\usepackage{fuzz}
\usepackage[colorlinks=true,linkcolor=blue,citecolor=blue,urlcolor=blue]{hyperref}

\begin{document}

\title{lux Agent-Menu Lifecycle: a Z Specification}
\author{Formal model of Hub agent-menu ownership, inbox delivery, and the
  session-departure cascade -- proving no false-positive delivery and no
  orphaned-owner menu across every interleaving}
\date{September 2026}
\maketitle

% ============================================================
\section{Introduction}
% ============================================================

An agent registers a menu bar with the MCP \texttt{menu\_set} tool; the Hub
records the bar under the registering session and, when the user clicks an
item, delivers a reserved \texttt{lux.menu} event to that session's inbox,
which the agent drains with \texttt{recv()}. Three pieces of state travel
together through this feature -- who owns which bar, whose inbox is live to
receive a click, and what happens to both when a session departs -- and the
defects PR~\#495's review found are all in the seams between them.

This model has one job: to state, over a bounded carrier and every
interleaving ProB can reach, the two safety properties the feature must
hold, and to show each of the three reported defects violates exactly one
of them.

\begin{description}
  \item[M1 --- no false-positive delivery.] A menu event is delivered to a
    live, drainable inbox, or the caller learns it was not delivered
    (\texttt{provider\_gone}). It is never reported delivered into a queue no
    \texttt{recv()} will drain. This is what D1 breaks: \texttt{offer} reads
    the queue under its lock, releases the lock, then puts -- so a
    \texttt{drop\_session} in the window lands the put in an orphan and
    \texttt{offer} still returns \texttt{True}. M1 is the \emph{orphan-put}
    property: it forbids a \texttt{delivered} report for a put into a queue that
    was \emph{already gone at put time}. The distinct question of a put that
    lands on a \emph{live} inbox which a departure then tears down --- the
    offer-then-drop-discards ordering a later review raised --- is a different
    property, treated in the companion refinement
    \texttt{menu\_lifecycle\_delivery.tex} (see
    Section~\ref{sec:delivery-refinement}).
  \item[MO --- menu-ownership integrity.] No session owns a menu bar in the
    registry unless it is live. Equivalently, \texttt{menuOwner} $\subseteq$
    \texttt{registered} on every reachable state. This is what D2 and D3
    break, from two different sides: D2, a departure that removes the session
    from the registry without pruning its bar (on the lease-lapse path);
    D3, a store that records a bar for a session already departed (the
    admit/store window).
\end{description}

M1 and MO are deliberately kept out of the state schema. As in
\texttt{connection\_lease\_reaping.tex}, a safety property placed in the
state invariant is enforced as a guard on every successor, silently
disabling the very operation that would break it -- so a violating
operation would become unreachable rather than a counterexample. The schema
below carries only structural coherence; M1 and MO are checked as the
reachability of their negations (Section~\ref{sec:verify}).

The model refines \texttt{connection\_lease\_reaping.tex}'s departure
cascade. That specification proved, for the client registry and scene
ownership, that leaving the registry and releasing what one owned are a
single event no matter which trigger fired it (its I5), extended to the
subscription/writer/inbox cascade legs (its I7). MO is that same shape
applied to a fourth cascade leg, the menu registry, and the fix D2 names --
binding \texttt{HubMenuRegistry.drop\_session} into the shared
\texttt{DepartureCascade} -- is exactly ``add one more sink that every
departure path already runs.''

Eight operations realize the fixed design. \texttt{Connect},
\texttt{Identify}, \texttt{Arm}, and \texttt{LeaseLapse} move a session
through its registration and lease states. \texttt{SetMenu} models
\texttt{MenuOperations.set\_menu} \emph{once fixed}: the admit-liveness
check and the bar store are one atomic step -- the D3 fix, holding
\texttt{StoreLock} across both. \texttt{SetMenuRefused} is its anonymous
counterpart. \texttt{DepartGraceful} and \texttt{DepartLapsed} are the two
departure triggers; both sweep the menu registry and the inbox as cascade
legs -- the D2 fix, one cascade every trigger runs. \texttt{MenuDeliver}
and \texttt{MenuProviderGone} model \texttt{MenuEventRouter.deliver}
\emph{once fixed}: the queue read and the put are one atomic step -- the D1
fix, \texttt{offer}'s get and put under one \texttt{\_inboxes\_lock} hold.

% ============================================================
\section{Basic Types}
% ============================================================

\begin{zed}
  [CONN]
\end{zed}

A connection identity -- the \texttt{ConnectionId} keying the client
registry, the inbox table, and the menu registry alike. Two \texttt{CONN}
values suffice to exhibit an owner departing while a delivery or a store is
mid-flight beside a second, unrelated session; the carrier is bounded by
ProB's \texttt{DEFAULT\_SETSIZE}.

% ============================================================
\section{Free Types}
% ============================================================

\begin{zed}
  DSTATE ::= dnone | ddelivered | dgone
\end{zed}

The outcome \texttt{MenuEventRouter.deliver} last reported to its caller:
nothing yet, \texttt{delivered}, or \texttt{provider\_gone}. Prefixed
constructors, never \fuzz's logical \texttt{true}/\texttt{false}, so a reader
is never tempted to read \texttt{reported = ddelivered} as a bare
proposition.

\begin{zed}
  FSTATE ::= fclear | fset
\end{zed}

A two-value flag, matching \texttt{connection\_lease\_reaping.tex}'s
\texttt{FSTATE}. Here it records whether a click reported delivered actually
landed in an orphaned queue.

\begin{zed}
  APHASE ::= aidle | astore
\end{zed}

Whether a \texttt{set\_menu} is between its admit-liveness check and its bar
store. In the fixed model \texttt{SetMenu} is atomic, so this never leaves
\texttt{aidle}; the D3 fidelity control splits \texttt{SetMenu} and drives
it to \texttt{astore}, the window the split exposes.

\begin{zed}
  OPHASE ::= oidle | oput
\end{zed}

Whether an \texttt{offer} is between its queue read and its put. In the
fixed model \texttt{MenuDeliver} is atomic, so this never leaves
\texttt{oidle}; the D1 fidelity control splits \texttt{offer} and drives it
to \texttt{oput}.

% ============================================================
\section{State}
% ============================================================

One flat schema, structural coherence only. The two safety properties (M1,
MO) are deliberately absent, so a violating operation is a counterexample
rather than a disabled transition.

\begin{schema}{MenuReg}
  registered : \power CONN \\
  identified : \power CONN \\
  lapsed : \power CONN \\
  hasInbox : \power CONN \\
  menuOwner : \power CONN \\
  reported : DSTATE \\
  lost : FSTATE \\
  admitPhase : APHASE \\
  admitFor : \power CONN \\
  offerPhase : OPHASE \\
  offerFor : \power CONN
\where
  identified \subseteq registered \\
  lapsed \subseteq registered \\
  \# admitFor \leq 1 \\
  \# offerFor \leq 1
\end{schema}

\texttt{registered} is the client registry's live set. \texttt{identified}
is the subset that has declared an identity -- the ownership gate
\texttt{menu\_set} demands; a session cannot be identified without being
registered, hence \texttt{identified} $\subseteq$ \texttt{registered}.
\texttt{lapsed} is the subset whose lease has expired; again a subset of
\texttt{registered}, since a departed connection is no longer a lease the
registry tracks. \texttt{hasInbox} is the set with a live, drainable inbox
queue (an armed writer, not yet dropped). \texttt{menuOwner} is the set with
a stored bar in \texttt{HubMenuRegistry.\_by\_owner}.

Crucially, \texttt{menuOwner} is \emph{not} constrained $\subseteq$
\texttt{registered} here: that containment is MO, the property this model
proves, not a structural fact assumed. \texttt{hasInbox} is likewise
unconstrained, so the D1 control can capture a live inbox and then have it
dropped from under the offer.

\texttt{reported} and \texttt{lost} carry the delivery outcome of the most
recent \texttt{deliver} call and whether any click has been lost to an
orphan. \texttt{admitPhase}/\texttt{admitFor} and
\texttt{offerPhase}/\texttt{offerFor} are the two split-window devices,
each a phase plus the at-most-one connection currently inside that window --
the same single-owner device \texttt{connection\_lease\_reaping.tex} uses
for \texttt{cascadeFor}, and inert (\texttt{aidle}/\texttt{oidle}, both sets
empty) throughout the fixed model.

% ============================================================
\section{Initialization}
% ============================================================

\begin{schema}{Init}
  MenuReg'
\where
  registered' = \emptyset \\
  identified' = \emptyset \\
  lapsed' = \emptyset \\
  hasInbox' = \emptyset \\
  menuOwner' = \emptyset \\
  reported' = dnone \\
  lost' = fclear \\
  admitPhase' = aidle \\
  admitFor' = \emptyset \\
  offerPhase' = oidle \\
  offerFor' = \emptyset
\end{schema}

Nothing has connected: every set empty, no delivery reported, nothing lost,
and neither split window open.

% ============================================================
\section{Operations}
% ============================================================

\subsection{Connect --- arrival and reconnect}

Models \texttt{HubDisplay.register\_client}: a session joins the registry
with a live (un-lapsed) lease. As in \texttt{connection\_lease\_reaping.tex}
there is no precondition privileging a fresh connection over a reconnecting
one, and nothing here touches another session's state. A reconnect does not
resubscribe or re-own anything -- \texttt{identified}, \texttt{hasInbox},
and \texttt{menuOwner} are left as departure left them.

\begin{schema}{Connect}
  \Delta MenuReg \\
  c? : CONN
\where
  registered' = registered \cup \{ c? \} \\
  lapsed' = lapsed \setminus \{ c? \} \\
  identified' = identified \\
  hasInbox' = hasInbox \\
  menuOwner' = menuOwner \\
  reported' = reported \\
  lost' = lost \\
  admitPhase' = admitPhase \\
  admitFor' = admitFor \\
  offerPhase' = offerPhase \\
  offerFor' = offerFor
\end{schema}

\subsection{Identify --- the ownership gate}

Models a session declaring an identity, the precondition
\texttt{MenuArming.admit} enforces. Only a registered session may identify.

\begin{schema}{Identify}
  \Delta MenuReg \\
  c? : CONN
\where
  c? \in registered \\
  identified' = identified \cup \{ c? \} \\
  registered' = registered \\
  lapsed' = lapsed \\
  hasInbox' = hasInbox \\
  menuOwner' = menuOwner \\
  reported' = reported \\
  lost' = lost \\
  admitPhase' = admitPhase \\
  admitFor' = admitFor \\
  offerPhase' = offerPhase \\
  offerFor' = offerFor
\end{schema}

\subsection{Arm --- bind the inbox writer}

Models \texttt{inbox.ensure\_writer}: a registered session gains a live
inbox queue. \texttt{SetMenu} arms the inbox as part of admitting the owner;
\texttt{Arm} is the standalone form (a plain \texttt{topic\_subscribe} also
reaches it), included so a delivery target's inbox can exist independently of
a menu store.

\begin{schema}{Arm}
  \Delta MenuReg \\
  c? : CONN
\where
  c? \in registered \\
  hasInbox' = hasInbox \cup \{ c? \} \\
  registered' = registered \\
  identified' = identified \\
  lapsed' = lapsed \\
  menuOwner' = menuOwner \\
  reported' = reported \\
  lost' = lost \\
  admitPhase' = admitPhase \\
  admitFor' = admitFor \\
  offerPhase' = offerPhase \\
  offerFor' = offerFor
\end{schema}

\subsection{LeaseLapse --- time passes with nothing to renew it}

Models \texttt{SessionLease.is\_live} going false: a registered lease
lapses. It gates \texttt{DepartLapsed} below, and it is the trigger
\texttt{connection\_lease\_reaping.tex} showed a purely pull-based design
cannot service on an idle Hub.

\begin{schema}{LeaseLapse}
  \Delta MenuReg \\
  c? : CONN
\where
  c? \in registered \\
  c? \notin lapsed \\
  lapsed' = lapsed \cup \{ c? \} \\
  registered' = registered \\
  identified' = identified \\
  hasInbox' = hasInbox \\
  menuOwner' = menuOwner \\
  reported' = reported \\
  lost' = lost \\
  admitPhase' = admitPhase \\
  admitFor' = admitFor \\
  offerPhase' = offerPhase \\
  offerFor' = offerFor
\end{schema}

\subsection{SetMenu --- admit and store, one atomic step (the D3 fix)}
\label{sec:setmenu}

Models \texttt{MenuOperations.set\_menu} for an admitted (registered and
identified) caller. \textbf{This is the D3 fix.} The shipped code calls
\texttt{MenuArming.admit} -- which renews the caller and arms its inbox under
\texttt{StoreLock}, then releases the lock -- and only afterward calls
\texttt{HubMenuRegistry.set\_menus} under the registry's own separate lock.
Here the renewal, the arming, and the store are one atomic step: no
departure can interleave between the liveness check and the store, which is
exactly what holding \texttt{StoreLock} across both buys.

\begin{schema}{SetMenu}
  \Delta MenuReg \\
  c? : CONN
\where
  c? \in registered \\
  c? \in identified \\
  lapsed' = lapsed \setminus \{ c? \} \\
  hasInbox' = hasInbox \cup \{ c? \} \\
  menuOwner' = menuOwner \cup \{ c? \} \\
  registered' = registered \\
  identified' = identified \\
  reported' = reported \\
  lost' = lost \\
  admitPhase' = admitPhase \\
  admitFor' = admitFor \\
  offerPhase' = offerPhase \\
  offerFor' = offerFor
\end{schema}

\subsection{SetMenuRefused --- anonymous ownership refused}

Models \texttt{admit} returning \texttt{False}: a caller that is not
identified (registered but anonymous, or not registered at all) is refused.
\texttt{renew\_if\_registered} still renews a registered caller's lease, but
nothing is armed and no bar is stored -- ``nothing anonymous owns a menu
item.''

\begin{schema}{SetMenuRefused}
  \Delta MenuReg \\
  c? : CONN
\where
  c? \notin identified \\
  lapsed' = lapsed \setminus \{ c? \} \\
  registered' = registered \\
  identified' = identified \\
  hasInbox' = hasInbox \\
  menuOwner' = menuOwner \\
  reported' = reported \\
  lost' = lost \\
  admitPhase' = admitPhase \\
  admitFor' = admitFor \\
  offerPhase' = offerPhase \\
  offerFor' = offerFor
\end{schema}

\subsection{DepartGraceful --- clean disconnect / SDK idle-reap}
\label{sec:departgraceful}

Models \texttt{HubDisplay.drop\_connection} together with
\texttt{session\_cleanup.py}'s teardown legs: the session leaves the
registry and every cascade leg it owns is released in the same step -- scene
ownership (elided here), its inbox, and its menu bar. Removal and release
are one predicate, so there is no state in which \texttt{c?} has left
\texttt{registered} under this operation without \texttt{menuOwner} and
\texttt{hasInbox} already reflecting the release.

\begin{schema}{DepartGraceful}
  \Delta MenuReg \\
  c? : CONN
\where
  c? \in registered \\
  registered' = registered \setminus \{ c? \} \\
  identified' = identified \setminus \{ c? \} \\
  lapsed' = lapsed \setminus \{ c? \} \\
  hasInbox' = hasInbox \setminus \{ c? \} \\
  menuOwner' = menuOwner \setminus \{ c? \} \\
  reported' = reported \\
  lost' = lost \\
  admitPhase' = admitPhase \\
  admitFor' = admitFor \\
  offerPhase' = offerPhase \\
  offerFor' = offerFor
\end{schema}

\subsection{DepartLapsed --- the timed backstop (the D2 fix)}
\label{sec:departlapsed}

Models \texttt{HubDisplay.reap\_lapsed\_leases} $\to$
\texttt{\_depart\_lapsed}: the background timer departs a lapsed connection
with no read, no other connection's write, and no disconnect signal
required. \textbf{This is the D2 fix.} In the shipped code this path runs
the \texttt{DepartureCascade}, which fires the inbox sink but not the menu
registry -- the menu prune is wired only to \texttt{session\_cleanup.py},
which the timer never runs. Here \texttt{DepartLapsed} sweeps
\texttt{menuOwner} in the very same predicate that sweeps \texttt{hasInbox}
and \texttt{registered}: the menu prune is one more cascade leg every
departure trigger runs, exactly as the inbox prune already is. Its predicate
is \texttt{DepartGraceful}'s with one added guard, \texttt{c? $\in$ lapsed}.

\begin{schema}{DepartLapsed}
  \Delta MenuReg \\
  c? : CONN
\where
  c? \in registered \\
  c? \in lapsed \\
  registered' = registered \setminus \{ c? \} \\
  identified' = identified \setminus \{ c? \} \\
  lapsed' = lapsed \setminus \{ c? \} \\
  hasInbox' = hasInbox \setminus \{ c? \} \\
  menuOwner' = menuOwner \setminus \{ c? \} \\
  reported' = reported \\
  lost' = lost \\
  admitPhase' = admitPhase \\
  admitFor' = admitFor \\
  offerPhase' = offerPhase \\
  offerFor' = offerFor
\end{schema}

\subsection{MenuDeliver --- read and put, one atomic step (the D1 fix)}
\label{sec:menudeliver}

Models \texttt{MenuEventRouter.deliver} landing a click on a live owner's
inbox. \textbf{This is the D1 fix.} The router reads the live set, then
calls \texttt{inbox.offer}, whose \texttt{get}-then-\texttt{put} the shipped
code splits across a lock release. Here the whole delivery -- the target is
live, its inbox is present, the message is put -- is one atomic step,
matching \texttt{offer} holding \texttt{\_inboxes\_lock} across its
\texttt{get} and its \texttt{put}. A live target with a live inbox is
reported \texttt{delivered}; nothing is ever put into an orphan, so
\texttt{lost} is untouched.

\begin{schema}{MenuDeliver}
  \Delta MenuReg \\
  c? : CONN
\where
  c? \in registered \\
  c? \in hasInbox \\
  reported' = ddelivered \\
  registered' = registered \\
  identified' = identified \\
  lapsed' = lapsed \\
  hasInbox' = hasInbox \\
  menuOwner' = menuOwner \\
  lost' = lost \\
  admitPhase' = admitPhase \\
  admitFor' = admitFor \\
  offerPhase' = offerPhase \\
  offerFor' = offerFor
\end{schema}

\subsection{MenuProviderGone --- the honest non-delivery}
\label{sec:providergone}

Models \texttt{deliver} returning \texttt{provider\_gone}: either the
router's live-set read fails (\texttt{c? $\notin$ registered}) or
\texttt{offer} finds no inbox to put into (\texttt{c? $\notin$ hasInbox}).
Either way the caller is told the truth -- not delivered -- and can re-push.
The negated conjunction is exactly the complement of \texttt{MenuDeliver}'s
guard, so between them every target is either delivered to or told
\texttt{provider\_gone}: there is no third outcome, and in particular no
silent drop.

\begin{schema}{MenuProviderGone}
  \Delta MenuReg \\
  c? : CONN
\where
  \lnot (c? \in registered \land c? \in hasInbox) \\
  reported' = dgone \\
  registered' = registered \\
  identified' = identified \\
  lapsed' = lapsed \\
  hasInbox' = hasInbox \\
  menuOwner' = menuOwner \\
  lost' = lost \\
  admitPhase' = admitPhase \\
  admitFor' = admitFor \\
  offerPhase' = offerPhase \\
  offerFor' = offerFor
\end{schema}

% ============================================================
\section{Invariants}
\label{sec:inv}
% ============================================================

Neither property is placed in \texttt{MenuReg}, for the reason
\texttt{connection\_lease\_reaping.tex} gives: a predicate there is enforced
as a guard on every successor, silently disabling the operation that would
break it.

\paragraph{M1 --- no false-positive delivery.}
\[
  \lnot\, (reported = ddelivered \land lost = fset).
\]
If the caller was told a click was delivered, that click is not sitting in
an orphaned queue. In the fixed spec no operation ever sets \texttt{lost} to
\texttt{fset}: \texttt{MenuDeliver} puts only onto a live inbox
(\texttt{hasInbox}) in the same atomic step it checks the target live, and
\texttt{MenuProviderGone} reports \texttt{dgone}, never \texttt{ddelivered}.
So the negation is not reachable. The D1 control
(Section~\ref{sec:fidelity-d1}) splits the delivery and shows the negation
becomes reachable the moment a \texttt{drop\_session} lands between the read
and the put.

\paragraph{MO --- menu-ownership integrity.}
\[
  \lnot\, \exists\, c : CONN @ c \in menuOwner \land c \notin registered.
\]
No session owns a menu bar unless it is live. Every operation that removes a
connection from \texttt{registered} -- \texttt{DepartGraceful} and
\texttt{DepartLapsed} -- removes it from \texttt{menuOwner} in the same
predicate (Sections~\ref{sec:departgraceful},~\ref{sec:departlapsed}); and
the one operation that adds to \texttt{menuOwner}, \texttt{SetMenu}, is
atomic and guarded on \texttt{c? $\in$ registered}, so it never records a bar
for a departed owner (Section~\ref{sec:setmenu}). MO is the general shape
both D2 and D3 break: D2, a departure trigger that removes from
\texttt{registered} without pruning \texttt{menuOwner}; D3, a store that
adds to \texttt{menuOwner} after the owner has left \texttt{registered}. A
single invariant rules out both, so a third departure trigger or a third
store path added later is checked against the same property, not a fresh ad
hoc patch -- exactly as \texttt{connection\_lease\_reaping.tex}'s I5 absorbs
its own fourth removal path.

\paragraph{The lock order MO's D3 half requires, stated explicitly.}
\label{sec:lockorder}
The design's Section~6 tripwire asked whether the fix introduces a new
lock-acquisition order. For M1 and D2 the answer is no. M1's fix is
atomicity \emph{within} \texttt{\_inboxes\_lock} alone -- \texttt{offer}'s
\texttt{get} and \texttt{put} under one hold of the lock it already takes --
adding no nesting; and \texttt{MenuEventRouter.deliver}'s live-set read
(client-registry lock, released) precedes the \texttt{offer}
(\texttt{\_inboxes\_lock}) with no nesting, the same order
\texttt{CallbackRouter.route} already uses. D2's fix binds one more sink
into the \texttt{DepartureCascade}, which already runs under
\texttt{StoreLock} on every departure path.

D3 is the exception the tripwire named, and it must be stated plainly rather
than waved past. The D3 fix requires the admit-liveness check and the
\texttt{set\_menus} store to fall under one hold of \texttt{StoreLock} -- so
the menu registry's own lock is acquired \emph{inside} \texttt{StoreLock},
a nesting the shipped code does not have (\texttt{set\_menus} runs under the
menu lock alone, after \texttt{StoreLock} has been released). This is a new
order: \texttt{StoreLock} $\to$ menu-registry lock. It introduces no
deadlock cycle, because it points the same direction as D2's fix: the
departure cascade's menu sink also fires under \texttt{StoreLock} and then
takes the menu lock (\texttt{StoreLock} $\to$ menu-registry lock). Both the
store path and the departure path acquire the two locks in the same order,
so the acquisition graph gains only edges pointing one way -- and a cycle
needs two. \texttt{StoreLock} is the outer lock on every path in the Hub
today; the menu registry lock joining as a consistently-inner lock is one
more caller taking the locks in the order every existing caller already
does. In this model that single-critical-section atomicity is represented by
\texttt{SetMenu} being one schema; the D3 control splits it and shows the
window the split exposes is where the orphaned-owner menu appears.

% ============================================================
\section{Fidelity Controls}
\label{sec:fidelity}
% ============================================================

Each control differs from this specification in exactly one place, and each
makes the corresponding invariant's negation \emph{reachable} where the
fixed spec makes it unreachable. A model that cannot reproduce the defect it
guards against is too abstract to trust; these three show it can.

\paragraph{D1 --- \texttt{menu\_lifecycle\_d1\_buggy.tex}.}
\label{sec:fidelity-d1}
\texttt{MenuDeliver} is replaced by \texttt{OfferBegin}/\texttt{OfferEnd},
splitting \texttt{offer} at the lock release the shipped code has between its
\texttt{get} and its \texttt{put}. \texttt{OfferBegin(c?)}, for a live
\texttt{c?} whose inbox is present, captures \texttt{c?} into
\texttt{offerFor} and enters \texttt{oput} -- it has snapshotted a queue
reference and will put into it. \texttt{OfferEnd} puts and reports
\texttt{ddelivered} \emph{unconditionally}; if \texttt{c?} has since lost its
inbox, the put lands in an orphan and \texttt{lost} becomes \texttt{fset}.
The trace \texttt{OfferBegin(c?)}\,;\,\texttt{DepartGraceful(c?)}\,;
\,\texttt{OfferEnd} reaches \texttt{reported = ddelivered $\land$ lost =
fset} -- M1's negation, \emph{FOUND}. The fix (atomic \texttt{MenuDeliver})
removes the window: no \texttt{Depart} can interleave a \texttt{get} and a
\texttt{put} that are one step.

\paragraph{D2 --- \texttt{menu\_lifecycle\_d2\_buggy.tex}.}
\label{sec:fidelity-d2}
\texttt{DepartLapsed} drops the \texttt{menuOwner' = menuOwner $\setminus$
\{c?\}} conjunct -- the timer's cascade reaching the inbox and the registry
but not the menu registry, exactly the shipped wiring. \texttt{SetMenu}
stays atomic (D3 fixed), isolating the defect to the departure path. The
trace \texttt{Connect(c?)}\,;\,\texttt{Identify(c?)}\,;\,\texttt{SetMenu(c?)}
\,;\,\texttt{LeaseLapse(c?)}\,;\,\texttt{DepartLapsed(c?)} reaches
\texttt{c? $\in$ menuOwner $\land$ c? $\notin$ registered} -- MO's negation,
\emph{FOUND}. The fix sweeps \texttt{menuOwner} on the lapsed path too.

\paragraph{D3 --- \texttt{menu\_lifecycle\_d3\_buggy.tex}.}
\label{sec:fidelity-d3}
\texttt{SetMenu} is replaced by \texttt{SetMenuAdmit}/\texttt{SetMenuStore},
splitting \texttt{set\_menu} at the \texttt{StoreLock} release between
\texttt{admit} and \texttt{set\_menus}. \texttt{SetMenuAdmit(c?)}, for an
admitted \texttt{c?}, renews and arms and enters \texttt{astore};
\texttt{SetMenuStore} records the bar \emph{unconditionally}, without
re-checking liveness. \texttt{DepartLapsed} stays a full sweep (D2 fixed),
isolating the defect to the store path. The trace
\texttt{Connect(c?)}\,;\,\texttt{Identify(c?)}\,;\,\texttt{SetMenuAdmit(c?)}
\,;\,\texttt{DepartGraceful(c?)}\,;\,\texttt{SetMenuStore} reaches
\texttt{c? $\in$ menuOwner $\land$ c? $\notin$ registered} -- MO's negation,
\emph{FOUND}: the departure swept a registry that did not yet hold the bar,
and the store then wrote it into a registry the owner has left. The fix
(atomic \texttt{SetMenu}, the \texttt{StoreLock} hold of
Section~\ref{sec:lockorder}) removes the window.

% ============================================================
\section{The Delivery Refinement: offer-then-drop-discards}
\label{sec:delivery-refinement}
% ============================================================

M1 above is the orphan-put property: no \texttt{delivered} report for a put
into a queue already gone at put time, closed by \texttt{offer} holding
\texttt{\_inboxes\_lock} across its \texttt{get} and \texttt{put}. A round-6
review of PR~\#495 raised a \emph{different} window, one this model's atomic
\texttt{MenuDeliver} abstracts away rather than rules on: the router's
live-set read passes, a departure then removes the connection, \texttt{offer}
then wins its lock and puts onto the \emph{still-present} inbox (returning
\texttt{True}), and only then the cascade's \texttt{drop\_session} discards
that inbox with the just-enqueued event in it. This is offer-\emph{then}-drop,
\emph{after} a fully-completed \texttt{offer} --- the D1 fix does not touch it,
because no lock can stop a recipient from departing once delivered to.

The abstraction is real: \texttt{MenuDeliver} here reports \texttt{ddelivered}
and carries no token for the enqueued message, and departure is one atomic
step, so this model cannot represent an enqueued event being discarded and
therefore cannot say whether that discard is a false delivery. The companion
\texttt{menu\_lifecycle\_delivery.tex} supplies exactly the missing detail: a
\texttt{pending} token for the undrained event, a \texttt{Drain} for
\texttt{recv()}, and a \emph{two-phase} departure whose registry-discard and
inbox-teardown are separate steps --- the ``separate critical sections'' the
review named. Over that refinement it proves

\[ \textbf{M1b} \quad \lnot\, (lostLive = fset), \]

no undrained event is discarded from a still-registered connection. The
discard can only ever strike a connection that has \emph{already} left the
registry, because the code's departure removes the connection
(\texttt{HubDisplay.\_depart} $\to$ \texttt{\_clients.discard}) before the
cascade fires \texttt{inbox.drop\_session}, under one \texttt{StoreLock} hold
--- the order the companion's \texttt{DepartRegistry}-before-\texttt{DepartInbox}
encodes and its reorder control shows to be load-bearing. A departed recipient
has no \texttt{recv()} to miss, so the loss is the correct disposal of a click
for a session that is gone, not a false delivery to a live one. The review's
further worry, that the click path ``will not re-push the menu,'' does not
follow: the same departure fires \texttt{MenuDepartureSink}
(\texttt{operations/menus.py}), which prunes the owner's bar and calls
\texttt{mark\_menus()} on every trigger, so the re-push happens on the
departure path regardless of what the click path reported.

% ============================================================
\section{Verification}
\label{sec:verify}
% ============================================================

Type-checking is a \texttt{fuzz} pass:
\begin{verbatim}
  fuzz docs/menu_lifecycle.tex
\end{verbatim}

M1 and MO are safety properties, checked by asking ProB whether their
negation is \emph{reachable}. A goal reported \emph{not found} means the
invariant holds on every reachable state.

\begin{verbatim}
  # M1 (goal NOT found in this, the fixed, spec)
  probcli docs/menu_lifecycle.tex -model_check -nodead \
    -goal "reported = ddelivered & lost = fset" \
    -p DEFAULT_SETSIZE 3

  # MO (goal NOT found in this, the fixed, spec)
  probcli docs/menu_lifecycle.tex -model_check -nodead \
    -goal "#c.(c : CONN & c : menuOwner & not(c : registered))" \
    -p DEFAULT_SETSIZE 3
\end{verbatim}

The positive outcomes are confirmed reachable, so the invariants are not
vacuously true of a model that can never deliver or store at all:

\begin{verbatim}
  # A click is delivered to a live inbox (goal FOUND)
  probcli docs/menu_lifecycle.tex -model_check -nodead \
    -goal "reported = ddelivered & lost = fclear" \
    -p DEFAULT_SETSIZE 3

  # A live, identified session owns a stored bar (goal FOUND)
  probcli docs/menu_lifecycle.tex -model_check -nodead \
    -goal "#c.(c : CONN & c : menuOwner & c : registered)" \
    -p DEFAULT_SETSIZE 3

  # provider_gone is reachable (goal FOUND)
  probcli docs/menu_lifecycle.tex -model_check -nodead \
    -goal "reported = dgone" \
    -p DEFAULT_SETSIZE 3
\end{verbatim}

Deadlock-freedom and the structural state invariant are checked by a full
model-check over the reachable state space:

\begin{verbatim}
  probcli docs/menu_lifecycle.tex -model_check -p DEFAULT_SETSIZE 3
\end{verbatim}

The fidelity controls make the negations \emph{FOUND}:

\begin{verbatim}
  # D1: M1's negation FOUND against the D1 control, NOT found against fixed
  probcli docs/menu_lifecycle_d1_buggy.tex -model_check -nodead \
    -goal "reported = ddelivered & lost = fset" \
    -p DEFAULT_SETSIZE 3

  # D2: MO's negation FOUND against the D2 control, NOT found against fixed
  probcli docs/menu_lifecycle_d2_buggy.tex -model_check -nodead \
    -goal "#c.(c : CONN & c : menuOwner & not(c : registered))" \
    -p DEFAULT_SETSIZE 3

  # D3: MO's negation FOUND against the D3 control, NOT found against fixed
  probcli docs/menu_lifecycle_d3_buggy.tex -model_check -nodead \
    -goal "#c.(c : CONN & c : menuOwner & not(c : registered))" \
    -p DEFAULT_SETSIZE 3
\end{verbatim}

Re-run \texttt{fuzz} and these goal checks whenever \texttt{inbox.py}'s
\texttt{offer}/\texttt{drop\_session}, \texttt{menu\_registry.py}'s
\texttt{set\_menus}/\texttt{drop\_session}, \texttt{menu\_event.py}'s
\texttt{MenuEventRouter.deliver}, \texttt{operations/menus.py}'s
\texttt{set\_menu}, \texttt{operations/menu\_arming.py}'s \texttt{admit},
\texttt{hub\_display.py}'s \texttt{\_depart}/\texttt{\_depart\_lapsed}/
\texttt{drop\_connection}, \texttt{departure\_cascade.py}, or
\texttt{session\_cleanup.py} change.

\end{document}
